problem:  
Let \( (M,g) \cong (G/K,g) \) be a **homogeneous Ricci soliton**, with \( G \) acting transitively by isometries and \( K \) compact.  
Let \( \mathfrak g = \mathfrak k \oplus \mathfrak p \) be a reductive decomposition, and let \( D\in \mathrm{Der}(\mathfrak g) \) denote the soliton derivation appearing in the algebraic soliton equation  
\[
\mathrm{Ric}_g = c\, I + S(D|_{\mathfrak p}), \qquad c \in \mathbb{R},
\]
where \( S(A)=\tfrac12(A+A^t) \) is the symmetric part relative to \(g\).  
Assume that the restriction \( D|_{\mathfrak n} \) to the nilradical \( \mathfrak n \) of \( \mathfrak g \) is diagonalizable over \( \mathbb C \) with real eigenvalues.

Define \( \mathcal{M}_G \) as the manifold of \( G \)-invariant metrics on \( G/K \).  
The **linearized soliton operator** at \(g\) can be viewed (modulo diffeomorphisms and scaling) as
\[
L_g(h) := \frac{d}{dt}\Big|_{t=0} \big(\mathrm{Ric}_{g_t} - c_t g_t - S(D_{t}|_{\mathfrak p})\big),
\]
acting on \( h \in T_g \mathcal{M}_G \simeq  \operatorname{Sym}^2(\mathfrak p^*) \).

**Question:**  
Is the kernel of the linearized soliton operator \(L_g\), modulo infinitesimal isometries and scaling, **isomorphic to the space of symmetric derivations** \( S \in \mathrm{Der}(\mathfrak g) \) that **commute with \(D\)**?  
Equivalently, are all first-order homogeneous Ricci soliton deformations driven solely by commuting symmetric derivations, or can genuinely new infinitesimal soliton directions arise from noncommuting perturbations?

If noncommuting infinitesimal deformations exist, can they be organized into an algebraic quotient space
\[
\mathcal{H}^1_D(\mathfrak g)
:= 
\frac{\{\,\Phi \in \mathrm{Der}(\mathfrak g) \mid [D,\Phi] \in \mathrm{Der}(\mathfrak g)\}}
{\{\, [D,\Psi] \mid \Psi \in \mathrm{Der}(\mathfrak g) \,\}},
\]
interpretable as a **“\(D\)-centralizer cohomology”**, analogously controlling obstructions to rigidity?  
Does there exist any example of a homogeneous Ricci soliton with nontrivial \(\mathcal{H}^1_D(\mathfrak g)\)?

Why is it a "good" problem:  
This problem precisely formulates the **infinitesimal rigidity** question for homogeneous Ricci solitons in an analytically and algebraically coherent way. It seeks to determine whether the soliton condition’s linearization admits nontrivial deformations beyond those generated by derivations commuting with the soliton derivation \(D\). A positive rigidity answer would identify algebraic centralization as the complete mechanism governing homogeneous soliton uniqueness; a negative answer would uncover a yet-unknown cohomological structure controlling soliton moduli—bridging Ricci flow analysis, the deformation theory of Lie algebras, and the structure of derivation spaces.  
The problem is elegant yet profound: its statement involves only familiar algebraic-geometric data (\(G/K\), \(D\), centralizers), but its resolution would unify geometric analysis and Lie algebra cohomology at a foundational level, exemplifying a "narrow entry, deep consequence" research problem.